\documentclass{beamer}
\usetheme{Madrid}
\title{Measurement Error and Misclassification}
\author{}
\date{}
\begin{document}
\frame{\titlepage}
\begin{frame}{Measurement Error and Misclassification}
\begin{itemize}
\item 1
\end{itemize}
\end{frame}
\begin{frame}{Study design in the presence of misclassification}
\begin{itemize}
\item 2
\end{itemize}
\end{frame}
\begin{frame}{Main study/internal validation study design  }
\begin{itemize}
\item 3
\end{itemize}
\end{frame}
\begin{frame}{Main study/external validation study design  }
\begin{itemize}
\item 4
\end{itemize}
\end{frame}
\begin{frame}{Two design problems are most important}
\begin{itemize}
\item 5
\end{itemize}
\end{frame}
\begin{frame}{Note that                  can be rewritten as             , }
\begin{itemize}
\item 6
\end{itemize}
\end{frame}
\begin{frame}{Greenland S. (1988)}
\begin{itemize}
\item 7
\end{itemize}
\end{frame}
\begin{frame}{Slide}
\end{frame}
\begin{frame}{Greenland S. (1988)}
\begin{itemize}
\item 9
\end{itemize}
\end{frame}
\begin{frame}{Greenland S. (1988)}
\begin{itemize}
\item 10
\end{itemize}
\end{frame}
\begin{frame}{Table 4: Optimal proportion          for validation substudy as a function of the cost ratio         }
\begin{itemize}
\item 11
\end{itemize}
\begin{itemize}
\item What is the relative cost of the measurement,    , for your study?   \\ For NHS, it has been estimated at 500.
\end{itemize}
\end{frame}
\begin{frame}{12}
\begin{itemize}
\item Assumes case-control ratios,         and        , are 1;  all other parameters given by the SIDS data. \\  \\ In case-control studies, with relatively low cost ratios (medical records vs. self-report), the fully validated design is often optimal \\  \\ Differential vs. non-differential misclassification assumption has a big impact on the design.
\end{itemize}
\end{frame}
\begin{frame}{Reilly M. (1996)}
\begin{itemize}
\item 13
\end{itemize}
\end{frame}
\begin{frame}{More complex designs have been developed in which the }
\begin{itemize}
\item 14
\end{itemize}
\end{frame}
\begin{frame}{Buonaccorsi (1990)}
\begin{itemize}
\item 15
\end{itemize}
\end{frame}
\begin{frame}{Note that                                                                                            ,}
\begin{itemize}
\item 16
\end{itemize}
\end{frame}
\begin{frame}{For             , the optimal MS/IVS needs to be compared to the fully}
\begin{itemize}
\item 17
\end{itemize}
\end{frame}
\begin{frame}{Buonaccorsi observed the following:}
\begin{itemize}
\item 18
\end{itemize}
\end{frame}
\begin{frame}{Spiegelman and Gray (1991)}
\begin{itemize}
\item 19
\end{itemize}
\end{frame}
\begin{frame}{Spiegelman and Gray (1991)}
\begin{itemize}
\item Considered MS/IVS, MS/EVS and fully validated studies, based on MLE. \\ Criterion is to minimize cost, subject to a specified discriminatory power. That is, the optimal design satisfies the power requirements of two hypothesis testing setups: First                plays the role of the null and                 plays the role of the alternative; then                is the null and                is the alternative.
\end{itemize}
\begin{itemize}
\item 20
\end{itemize}
\end{frame}
\begin{frame}{Discriminatory power}
\begin{itemize}
\item Power to discriminate two parameter values. \\ Let                  and                be the two values that we want to discriminate, and let (e, f) be the 95\% CI for log(OR). Then, the power to discriminate               and                 is the smaller of   \\    and                                               .    \\ Since                                                 \\                                                    , it has one side 	type I error rate of 2.5\% for each direction. \\ Greenland S. On Sample-Size and Power Calculations for Studies Using Confidence Intervals. AJE 1988; 128, 231-235.
\end{itemize}
\begin{itemize}
\item 21
\end{itemize}
\end{frame}
\begin{frame}{Discriminatory power}
\begin{itemize}
\item (e                         f)  ----standard \\          (e                   f)                ----additional \\    \\ A sample size may be considered adequate for discriminating between two parameter values (one of which is correct), if it is large enough so that there is a good chance that the observed CI excludes the incorrect value.
\end{itemize}
\begin{itemize}
\item 22
\end{itemize}
\end{frame}
\begin{frame}{Spiegelman and Gray (1991)}
\begin{itemize}
\item 23
\end{itemize}
\begin{itemize}
\item It assumes measurement error model: \\  \\  \\ Recall that if                               , classical measurement error                                      \\     can be written as                                                 with
\end{itemize}
\end{frame}
\begin{frame}{24}
\end{frame}
\begin{frame}{In this figure,                                    and                         under classical measurement error model for MS/EVS design.}
\begin{itemize}
\item 25
\end{itemize}
\end{frame}
\begin{frame}{Rosner, Spiegelman and Willett (Rosner et al., 1992)}
\begin{itemize}
\item Rosner, Spiegelman and Willett (Rosner et al., 1992) investigated  \\ the role of        (size of reliability study) and       (number of replicates per participant in the reliability study)  on the efficiency of the regression calibration estimates for the effects of systolic blood pressure (SBP), serum cholesterol (CHOL), serum glucose (GLUC) and BMI on the 10-year cumulative  incidence of CHD, based on the Framingham Heart Study.  There were 1,731 men in the main study cohort. \\  \\  \\ Rosner B., Spiegelman D., Willett W.C. (1992) Correction of logistic regression relative risk \\ estimates and confidence intervals for random within-person measurement error. Am J Epidemiol  \\ 136:1400-13.
\end{itemize}
\begin{itemize}
\item 26
\end{itemize}
\end{frame}
\begin{frame}{27}
\begin{itemize}
\item It is apparent that as         increases,  \\ the variances of the measurement \\ error corrected log relative risks \\ using regression calibration  converge to an asymptote.  \\   \\  \\  \\  \\  \\ At small values of         , larger  \\ values of           improve the  \\ precision, but once         is greater  \\ than 100 or so, there  is very little  \\ advantage over                .
\end{itemize}
\end{frame}
\begin{frame}{Concepts}
\begin{itemize}
\item 28
\end{itemize}
\end{frame}
\end{document}